\scanpage{12}
\uline{Prop} \unsure{1|1a|1\textsuperscript{o}} Soit $f:X\to Y$ un morphisme \deleted{de type fini} \deleted{\unsure{univ\textsuperscript{t} ouvert|loc.}} dont les fibres sont \deleted{\unsure{ult.}} irréductibles, Pour tt $y\in Y$, soit $\varphi(y)$ le p\textsuperscript{t} générique de $X_y$. \deleted{Pour que} \inserted{Conditions équivalentes} (i) $\varphi$ \deleted{soit} continue, \deleted{il f. et s.} \deleted{\unsure{(ii)} $y\in\overline{y'}\Rightarrow\varphi(y)\in\overline{\varphi(y')}$} \deleted{que} (ii) $f$ \deleted{soit} \uline{ouvert}. Cela implique la condition suivante, \uline{et} lui est équivalente si $Y$ est loc. noeth., $f$ de type fini :

(iii) $y\in\overline{y'}\Rightarrow\varphi(y)\in\overline{\varphi(y')}$.
\deleted{\unsure{$\Longleftrightarrow$ (ii bis)}}

\uline{Dém} \deleted{On a évidemment (i) $\Rightarrow$ (ii).} \inserted{\deleted{En effet} On a \unsure{(i)|(ii)\textsuperscript{o}} $\Longleftrightarrow$ (ii), car \deleted{\unsure{soit}} $U$ \inserted{est} un ouvert de $X$, alors $\varphi^{-1}(U)=f(U)$. De plus (ii) $\Rightarrow$ (iii) trivial.} De plus (iii) équivaut à \inserted{(iii bis) Pour tt $y,y'$ tels que $y\in\overline{y'}$, et $x$ au dessus de $y$, \deleted{tel que} il existe $x'$ au dessus de $y'$ tel que $x\in\overline{x'}$.}

En effet, \deleted{on peut supposer} (ii) $\Rightarrow$ (ii bis), car on prend $x'=\varphi(y')$, et on a $x\in\overline{\varphi(y)}$ et $\varphi(y)\in\overline{\varphi(y')}$ d'où $x\in\overline{\varphi(y')}$. On a (ii bis) $\Rightarrow$ (ii), car \deleted{il existe} si $y\in\overline{y'}$, il existe $x'$ sur $y'$ tel que $\varphi(y)\in\overline{x'}$, et comme $x'\in\overline{\varphi(y')}$, on aura $\varphi(y)\in\overline{\varphi(y')}$.

Or on sait par ailleurs que \inserted{si $Y$ loc. nth, $f$ de type fini, alors} \deleted{\unsure{(ii bis)}} \inserted{(iii bis)} $\Rightarrow$ (ii) \deleted{\unsure{donc On}} (grâce à l'hyp. \unsure{neth.}) Donc (i) $\Rightarrow$ (ii) $\Longleftrightarrow$ (ii bis) $\Longleftrightarrow$ (iii)

\uline{Déf} 2 On dit alors que $f$ est \uline{\unsure{monique|monigue|monogène}}. Définition de Universellement monique.

\uline{Déf\textsuperscript{on}} 3 Une \deleted{morphisme application} \inserted{section \unsure{ensembliste}} $\varphi:Y\to X$ est dite \uline{section} \uline{\unsure{monisante|moniquante|monigante|monogène}} si il existe un ouvert $U$ de $X$, universellement monique sur $Y$, tel que $\varphi$ soit la fonction $\varphi_{U/Y}$.
% archive sheet number: 11
